\chapter{Topología y Entorno Tecnológico}

El punto de partida de toda investigación pericial informática exige una caracterización minuciosa del objetivo: su arquitectura de hardware, sistema operativo, topología de red y los servicios de red expuestos al momento del incidente. En el contexto de 2004, los servidores departamentales solían asumir roles múltiples y concurrentes (correo, archivos NFS, estaciones de cómputo y entornos gráficos X11), lo que ampliaba notablemente su superficie de ataque.

\begin{notaforense}[Contexto de Infraestructura en 2004]
El servidor estudiado representaba el estándar de cómputo de alto rendimiento departamental a principios de la década del 2000: una arquitectura x86 monoprocesador Intel Pentium III sobre bus PCI y discos duros SCSI/IDE, enlazada mediante conmutadores Fast Ethernet (100 Mbps) a una red universitaria con direccionamiento IPv4 público directo y sin cortafuegos de aplicación intermedios.
\end{notaforense}

\section{Descripción del sistema}

 \textbf{Nombre:} victima1

 \textbf{Dominio:} foo.bar.mx

 \textbf{Ip:} 192.0.2.10

 \textbf{Funciónes:} Servidor de correo, servidor de archivos, servidor gráfico.

 \textbf{Marca:} Genérico ensamblado.

 \textbf{Procesador:} Intel Pentium III 750 MHz.

 \textbf{Memoria:} 256 MB RAM.

 \textbf{Discos:} IDE ATA 66 20 GB.
 
 \textbf{Sistema Operativo:} GNU/Linux Debian 3.0-r1.  

\section{Ubicación física del sistema}

 El equipo estudiado estaba instalado en el edificio xxx del campus. Su conexión al switch general era por medio de una 
tarjeta de red ethernet con topología de estrella. Compartía directorios con el servidor \verb+xxx.xxx.xxx.2+ 
y había acceso para usuarios por consola y terminal gráfica.

\begin{perspectiva20}[Perspectiva 20 Años Después: Topología y Perímetro]
En 2004, era una práctica habitual conectar servidores de cómputo directamente con direcciones IP públicas sin capas de inspección profunda de paquetes ni arquitecturas de confianza cero (\textit{Zero Trust}). Actualmente, un sistema análogo operaría dentro de una VPC privada, aislado tras pasarelas seguras, balanceadores de carga con WAF, segmentación por microperímetros y autenticación reforzada mediante llaves criptográficas basadas en hardware.
\end{perspectiva20}
